% Shared packages and setup (common across Parts A/B)
\usepackage[T1]{fontenc}
\usepackage[utf8]{inputenc}
\usepackage{lmodern}
\usepackage{amsmath,amssymb,amsthm}
\usepackage{mathrsfs}
\usepackage{bm}
\usepackage{graphicx}
\usepackage{xcolor}
\usepackage{enumitem}
\usepackage{booktabs}
\usepackage{tabularx}
\usepackage{array}
\usepackage{listings}
\usepackage{microtype}
\usepackage{xstring} % string ops for robust const refs/labels
\usepackage{mdframed}
% Cross-document references (must precede hyperref for xr-hyper)
\usepackage{xr-hyper}
\makeatletter
% Prefer Wrapper B AUX when available; fallback to A. Relative to source dir.
\@ifundefined{IfFileExists}{}{%
  % Part I
  \IfFileExists{../../build/JSG_I_B.aux}{\externaldocument[JSGI-]{../../build/JSG_I_B}}{%
    \IfFileExists{../../build/JSG_I_A.aux}{\externaldocument[JSGI-]{../../build/JSG_I_A}}{}
  }%
  % Part II
  \IfFileExists{../../build/JSG_II_B.aux}{\externaldocument[JSGII-]{../../build/JSG_II_B}}{%
    \IfFileExists{../../build/JSG_II_A.aux}{\externaldocument[JSGII-]{../../build/JSG_II_A}}{}
  }%
}
\makeatother
\usepackage{hyperref}
\usepackage[backend=bibtex]{biblatex}
% --- Unicode fallbacks for occasional stray bytes in source files ---
% Map problematic Unicode to LaTeX-safe sequences to prevent inputenc errors
\DeclareUnicodeCharacter{009D}{}
\DeclareUnicodeCharacter{2013}{--}
\DeclareUnicodeCharacter{2014}{---}
\DeclareUnicodeCharacter{2018}{`}
\DeclareUnicodeCharacter{2019}{'}
\DeclareUnicodeCharacter{201C}{``}
\DeclareUnicodeCharacter{201D}{''}
% Hyperref setup
\hypersetup{
  colorlinks=true,
  linkcolor=blue!60!black,
  citecolor=blue!60!black,
  urlcolor=blue!60!black
}
% mdframed styles (rules / constants / ledger)
\definecolor{rulebg}{RGB}{248,248,250}
\definecolor{ruleframe}{RGB}{165,165,165}
\newmdenv[
  linecolor=ruleframe, linewidth=0.6pt, roundcorner=4pt,
  backgroundcolor=rulebg, innertopmargin=6pt, innerbottommargin=6pt,
  innerleftmargin=8pt, innerrightmargin=8pt, skipabove=8pt, skipbelow=8pt
]{rulebox}
\newenvironment{rulex}[1]{\begin{rulebox}\noindent\textbf{Rule (#1).}\ }{\end{rulebox}}
\definecolor{constbg}{RGB}{246,249,255}
\definecolor{constframe}{RGB}{140,160,190}
\newmdenv[
  linecolor=constframe, linewidth=0.6pt, roundcorner=4pt,
  backgroundcolor=constbg, innertopmargin=6pt, innerbottommargin=6pt,
  innerleftmargin=8pt, innerrightmargin=8pt, skipabove=10pt, skipbelow=10pt
]{constbox}
\definecolor{ledgerbg}{RGB}{255,251,235}
\definecolor{ledgerframe}{RGB}{230,200,120}
\newmdenv[
  linecolor=ledgerframe, linewidth=0.6pt, roundcorner=4pt,
  backgroundcolor=ledgerbg, innertopmargin=6pt, innerbottommargin=6pt,
  innerleftmargin=8pt, innerrightmargin=8pt, skipabove=8pt, skipbelow=8pt
]{ledgerbox}
% Listings style for API/code stubs
\definecolor{codebg}{RGB}{248,248,250}
\definecolor{codeframe}{RGB}{180,180,180}
\definecolor{codekw}{RGB}{0,0,160}
\definecolor{codestr}{RGB}{120,50,0}
\definecolor{codecm}{RGB}{90,90,90}
\lstdefinestyle{apicode}{
  basicstyle=\ttfamily\small,
  backgroundcolor=\color{codebg},
  frame=single, rulecolor=\color{codeframe}, frameround=tttt,
  columns=flexible, keepspaces=true, showstringspaces=false,
  keywordstyle=\color{codekw}\bfseries,
  stringstyle=\color{codestr}, commentstyle=\color{codecm}\itshape,
  tabsize=2, belowskip=0.6\baselineskip, aboveskip=0.6\baselineskip,
  numbersep=6pt, numbers=left, numberstyle=\tiny\color{codecm},
  language=Python
}
\lstset{style=apicode}
% Theorem styles
\theoremstyle{plain}
\newtheorem{theorem}{Theorem}
\newtheorem{lemma}{Lemma}
\newtheorem{proposition}{Proposition}
\newtheorem{corollary}{Corollary}
\theoremstyle{definition}
\newtheorem{definition}{Definition}
\newtheorem{assumption}{Assumption}
\theoremstyle{remark}
\newtheorem{remark}{Remark}
% Handy macros and shorthands
\newcommand{\R}{\mathbb{R}}
\newcommand{\E}{\mathbb{E}}
\newcommand{\FIM}{\mathrm{FIM}}
\newcommand{\normFIM}[1]{\left\|#1\right\|_{\FIM^{-1}}}
\newcommand{\uses}[1]{\par\smallskip\noindent\textit{Uses: }#1\par\smallskip}
\newenvironment{proofsketch}{\begin{proof}[Proof sketch]}{\end{proof}}
\newcommand{\absconst}[1]{c_{#1}}
\newcommand{\anglesep}{\theta_{\min}}
\newcommand{\cutq}{q}
\newcommand{\cont}{{\mathrm{cont}}}
\newcommand{\atom}{{\mathrm{atom}}}
\newcommand{\layertag}[1]{\textup{[#1]}}
% KPI & clamp shorthand
\newcommand{\RegretM}{\mathrm{Regret}_M}
\newcommand{\Hol}{\mathrm{Hol}}
\newcommand{\Res}{\mathrm{Res}}
\newcommand{\Lat}{\mathrm{Lat}}
\newcommand{\cgreen}{\chi_{\mathrm{green}}}
\newcommand{\calarm}{\chi_{\mathrm{alarm}}}
\newcommand{\cclear}{\chi_{\mathrm{clear}}}

% --- Paper-tenets helpers (Symbol Registry, Acceptance checks, Constants) ---
% Clickable constant labels (pair with \constref in macros.tex)
% Usage: place \constlabel{Smax} near the definition of S_max to anchor \constref{Smax}
\newcommand{\constlabel}[1]{\label{const:#1}}

% Symbol Registry environment and row helper
% Usage:
%   \begin{SymbolRegistry}
%     \symrow{C1}{$\kappa_{\max},\,\theta_{\min},\,r_{\mathrm{tube}}$}{Frontier geometry: curvature, incidence, tube radius.}
%     ...
%   \end{SymbolRegistry}
\newenvironment{SymbolRegistry}{%
  \begin{center}
  \renewcommand{\arraystretch}{1.2}%
  \begin{tabularx}{\textwidth}{@{} c >{\raggedright\arraybackslash}X >{\raggedright\arraybackslash}X @{} }
  \toprule
  \textbf{Code} & \textbf{Symbol(s)} & \textbf{Meaning / where used} \\
  \midrule
}{%
  \bottomrule
  \end{tabularx}
  \end{center}
}
\newcommand{\symrow}[3]{#1 & #2 & #3 \\ \addlinespace}

% Acceptance-checks box (ledger styling)
% Usage: \begin{acceptbox}[Acceptance checks for Section X] ... \end{acceptbox}
\newenvironment{acceptbox}[1][Acceptance checks]{%
  \begin{ledgerbox}\noindent\textbf{#1.}\ }{\end{ledgerbox}}

% Standing-constants box (alias of constbox with a title)
% Usage: \begin{standbox}[Standing constants] ... \end{standbox}
\newenvironment{standbox}[1][Standing constants]{%
  \begin{constbox}\noindent\textbf{#1.}\ }{\end{constbox}}

% Lightweight API code block alias (uses the global lstlisting style=apicode)
% Usage: \begin{apicodeblock} ... \end{apicodeblock}
\lstnewenvironment{apicodeblock}{}{}

% --- Tenets compatibility: constant refs (fallbacks if macros.tex not loaded) ---
% These provide the \constref and \constsym* macros used by the Symbol Registry.
% If macros.tex is also \input, its definitions will take precedence.
\providecommand{\constref}[1]{\hyperref[const:#1]{\ensuremath{\csname constsym#1\endcsname}}}
\providecommand{\constsymSmax}{S_{\max}}
\providecommand{\constsymGmax}{G_{\max}}
\providecommand{\constsymkappamax}{\kappa_{\max}}
\providecommand{\constsymthetamin}{\theta_{\min}}
\providecommand{\constsymLF}{L_{\mathrm F}}
\providecommand{\constsymrtube}{r_{\mathrm{tube}}}
\providecommand{\constsymtaumin}{\tau_{\min}}
\providecommand{\constsymLpsi}{L_\psi}
\providecommand{\constsymF}{F}
\providecommand{\constsymchi}{\chi}
\providecommand{\constsymepsilon}{\epsilon}
\providecommand{\constsymrho}{\rho}
\providecommand{\constsymtau}{\tau}
\providecommand{\constsymetastar}{\eta^\star}

% --- QOL: Dense lists for acceptance checks / KPIs ---------------------------
% Tight itemize/enumerate environments for audit-style lists
\newlist{tightitem}{itemize}{1}
\setlist[tightitem]{leftmargin=*, itemsep=2pt, topsep=2pt, parsep=0pt, partopsep=0pt}
\newlist{tightenum}{enumerate}{1}
\setlist[tightenum]{leftmargin=*, itemsep=2pt, topsep=2pt, parsep=0pt, partopsep=0pt, label=(\roman*)}

% --- QOL: Badges, statuses, and simple field printer ------------------------
\providecommand{\hash}[1]{\texttt{#1}}
\providecommand{\field}[2]{\noindent\textbf{#1:}\ #2}
\providecommand{\ok}{\textcolor{green!60!black}{\ensuremath{\checkmark}}}
\providecommand{\fail}{\textcolor{red!70!black}{\ensuremath{\times}}}
\providecommand{\accept}{\textcolor{green!60!black}{accept}}
\providecommand{\reject}{\textcolor{red!70!black}{reject}}

% Risk/clamp badges (deterministic replayable display)
\providecommand{\RiskGreen}{\colorbox{green!12}{\textsf{\footnotesize RISK: GREEN}}}
\providecommand{\RiskYellow}{\colorbox{yellow!20}{\textsf{\footnotesize RISK: YELLOW}}}
\providecommand{\RiskRed}{\colorbox{red!15}{\textsf{\footnotesize RISK: RED}}}

% IV Badge compact printer (values are free-form for operator control)
\providecommand{\IVbadge}[1]{\colorbox{blue!8}{\textsf{\footnotesize IV-Badge}}\ #1}

% --- QOL: LTL/Liveness operators and PR/CVaR operators ----------------------
\providecommand{\G}{\mathsf{G}}
\providecommand{\F}{\mathsf{F}}
\providecommand{\X}{\mathsf{X}}
\providecommand{\U}{\mathsf{U}}
\providecommand{\PR}{\mathrm{PR}}
\DeclareMathOperator{\CVaR}{CVaR}

% --- QOL: Confidence-interval, window, and paragraph head helpers -----------
\providecommand{\CI}[2]{\ensuremath{[\,#1,\ #2\,]}}
\providecommand{\win}[1]{\ensuremath{[t-#1,\ t]}}
\providecommand{\parhead}[1]{\par\noindent\textbf{#1.}\ }

% --- QOL: Witness/Audit boxes (mdframed presets) ---------------------------
\definecolor{witnessbg}{RGB}{243,248,255}
\definecolor{witnessframe}{RGB}{120,150,200}
\newmdenv[
  linecolor=witnessframe, linewidth=0.6pt, roundcorner=4pt,
  backgroundcolor=witnessbg, innertopmargin=6pt, innerbottommargin=6pt,
  innerleftmargin=8pt, innerrightmargin=8pt, skipabove=8pt, skipbelow=8pt
]{witnessbox}

\definecolor{auditbg}{RGB}{240,255,245}
\definecolor{auditframe}{RGB}{110,180,140}
\newmdenv[
  linecolor=auditframe, linewidth=0.6pt, roundcorner=4pt,
  backgroundcolor=auditbg, innertopmargin=6pt, innerbottommargin=6pt,
  innerleftmargin=8pt, innerrightmargin=8pt, skipabove=8pt, skipbelow=8pt
]{auditbox}

% --- QOL: Table column shortcuts (for compact registries/ledgers) -----------
\newcolumntype{L}[1]{>{\raggedright\arraybackslash}p{#1}}
\newcolumntype{C}[1]{>{\centering\arraybackslash}p{#1}}
\newcolumntype{R}[1]{>{\raggedleft\arraybackslash}p{#1}}
